\documentclass[11pt]{article}

\usepackage[a4paper,margin=0.88in]{geometry}
\usepackage{amsmath,amssymb,amsthm,mathtools}
\usepackage[hidelinks]{hyperref}
\usepackage{microtype}

\newtheorem{theorem}{Theorem}

\title{An Infinite Family of Integer Solutions to
\texorpdfstring{$z^2+y^2z+2x^3+1=0$}{z^2+y^2z+2x^3+1=0}}
\author{}
\date{}

\begin{document}
\maketitle
\vspace{-1.2em}

\begin{abstract}
We give an explicit polynomial family of integer solutions to
\(z^2+y^2z+2x^3+1=0\).  The proof is a direct substitution.  We then
briefly explain the sequence of ansatzes from which the family was found.
\end{abstract}

\section{The family}

\begin{theorem}\label{thm:main}
For every odd integer \(r\geq 3\), put
\[
\begin{aligned}
 x_r&=-\frac{r^8-8r^6+22r^4-22r^2+5}{2},\\
 y_r&=r^3-3r,\\
 z_r&=\frac{r^{12}-12r^{10}+57r^8-134r^6
              +159r^4-84r^2+11}{2}.
\end{aligned}
\]
Then \((x_r,y_r,z_r)\in\mathbb Z^3\) and
\[
        z_r^2+y_r^2z_r+2x_r^3+1=0.
\]
Consequently, the equation has infinitely many integer solutions.
\end{theorem}

\begin{proof}
Set
\[
 a=r^2-2,\qquad n=\frac{a^2-1}{2},\qquad
 T=\frac{r(a-1)}2,
\]
\[
 X=2n^2+a,\qquad M=4n^3+3an+1.
\]
Since \(r\) and \(a\) are odd, all five quantities are integers.  Direct
expansion gives
\[
 -X=x_r,\qquad 2T=y_r,\qquad M-2T^2=z_r,
\]
and
\[
 M^2-2X^3+1
 =\frac{(a-1)^4(a+2)^2}{4}
 =4\left(\frac{r(a-1)}2\right)^4
 =4T^4.
\]
Therefore, with \((x,y,z)=(-X,2T,M-2T^2)\),
\begin{align*}
 z^2+y^2z+2x^3+1
 &=z(z+y^2)-2X^3+1\\
 &=(M-2T^2)(M+2T^2)-2X^3+1\\
 &=M^2-2X^3+1-4T^4=0.
\end{align*}
Finally, \(y_r=r^3-3r\) is strictly increasing for \(r\geq3\), so the
solutions obtained from distinct odd values of \(r\) are distinct.
\end{proof}

\section{How the construction was found}\label{sec:discovery}

This section records the search that led to the formulas; it is not an
additional proof.

Write \(x=-X\).  The equation becomes
\[
        z(z+y^2)=2X^3-1.
\]
Thus the aim is to factor \(2X^3-1\) into two integers whose difference is
a square.  Taking \(y=2T\) and centering the factors symmetrically,
\[
        z=M-2T^2,\qquad z+y^2=M+2T^2,
\]
reduces the problem to
\begin{equation}\label{eq:target}
        M^2-2X^3+1=4T^4.
\end{equation}

There is an obvious polynomial solution of the unperturbed relation
\(M^2=2X^3\):
\[
        X=2n^2,\qquad M=4n^3.
\]
This suggests perturbing it by writing
\[
        X=2n^2+a,
        \qquad
        M=4n^3+bn+c.
\]
In the expansion of \(M^2-2X^3+1\), the coefficient of \(n^4\) is
\(8(b-3a)\).  Choosing \(b=3a\) removes this highest-order error, and the
simple choice \(c=1\) gives
\[
        X=2n^2+a,
        \qquad
        M=4n^3+3an+1,
\]
with residual
\[
        M^2-2X^3+1
        =8n^3-3a^2n^2+6an-2a^3+2.
\]

The leading terms indicate that \(n\) should be approximately \(a^2/2\).
Accordingly, one tries
\[
        n=\frac{a^2+pa+q}{2}
\]
and asks that four times the residual be the square of a monic cubic in
\(a\).  Elementary coefficient matching gives the particularly simple
choice
\[
        p=0,\qquad q=-1,
\]
so that
\[
        n=\frac{a^2-1}{2}
\]
and
\begin{align*}
 4(M^2-2X^3+1)
   &=(a^3-3a+2)^2\\
   &=(a-1)^4(a+2)^2.
\end{align*}
The factor \((a-1)^4\) is already a fourth power.  To make the remaining
factor a fourth power as well, impose
\[
        a+2=r^2,
        \qquad\text{that is,}\qquad
        a=r^2-2.
\]
Then
\[
        M^2-2X^3+1
        =4\left(\frac{r(a-1)}2\right)^4,
\]
which is exactly \eqref{eq:target} with
\(T=r(a-1)/2\).  Reversing the substitutions produces the family in
Theorem \ref{thm:main}.  The large polynomials in \(r\) arise only from
eliminating the intermediate variables \(a,n,T,X,M\).

\end{document}
